\section{\pimattn{}: Heterogeneous Decomposed Attention}
\label{sec:method}

\subsection{Overview}

\pimattn{} adopts a heterogeneous execution model where different stages of asymmetric quantized attention are mapped to their optimal compute substrate.
Figure~\ref{fig:overview} illustrates the overall architecture.

% TODO Figure 4: System overview — Q@K(decomposed)→PIM, Softmax→GPU, Score@V→GPU, with data flow arrows

The execution proceeds in three phases:
\begin{enumerate}
\item \textbf{Phase 1 (PIM)}: The decomposed Q@K dot product is computed on near-bank PIM units. The 2-bit KV cache data resides in PIM banks. Query $Q$ (float16) is broadcast to all banks. Each bank computes its portion of all four terms in Eq.~\ref{eq:full_expansion} using efficient 2-bit PEs.
\item \textbf{Phase 2 (GPU)}: Partial results from all banks are reduced (via simple summation). Softmax is computed on GPU using the aggregated attention scores.
\item \textbf{Phase 3 (GPU)}: The attention output is computed as $S \cdot V$, where $S$ is the float16 softmax output and $V$ is dequantized from 2-bit to float16 on GPU (leveraging GPU's massive float16 throughput).
\end{enumerate}

\subsection{Decomposed Q@K on Near-Bank PIM}

The core innovation of \pimattn{} is mapping the decomposed asymmetric dot product (Eq.~\ref{eq:full_expansion}) onto near-bank PIM.
For a query $Q \in \mathbb{R}^{M \times d}$ and key cache $K \in \mathbb{R}^{N \times d}$ (where $M=1$ in decode mode, $d$ is head dimension, $N$ is sequence length), each bank processes its partition of the data.

\textbf{Data distribution}: The KV cache is striped across $N_{\text{banks}}$ PIM banks.
Each bank holds a contiguous chunk of the $K$ and $V$ matrices along the sequence dimension.
Query $Q$ is broadcast to all banks (cost modeled as additional row-buffer fill time).

\textbf{Per-bank computation}: For its assigned chunk of $k_{\text{chunk}}$ key vectors, each bank computes:
\begin{equation}
\begin{aligned}
\text{(a)} &\quad \sum_{j} \hat{q}_j \hat{k}_{ij} \quad \text{(2-bit MAC)} \\
\text{(b)} &\quad \sum_{j} \hat{q}_j \quad \text{(2-bit accumulate, once per query)} \\
\text{(c)} &\quad \sum_{j} \hat{k}_{ij} \quad \text{(2-bit accumulate, per key vector)} \\
\text{(d)} &\quad \text{aggregate via Eq.~\ref{eq:full_expansion}}
\end{aligned}
\end{equation}

Operations (a)-(c) all operate on 2-bit data, making them highly efficient on simple PIM PEs.
The final float16 scaling by $s_Q s_K$ and subtraction of the constant term $d z_Q z_K$ are performed after reduction.

\subsection{PIM Latency Model}

We model the per-bank latency using a pipelined row-buffer execution model:

\begin{algorithm}[t]
\caption{Per-Bank PIM GEMV Latency}
\label{alg:pim}
\SetAlgoLined
\KwIn{GEMV dimensions $M,K,N$, element size $e$ bytes, bank config}
\KwOut{Latency breakdown (rowbuffer time, PE time)}

$B_{\text{total}} \gets (MK + KN + MN) \cdot e$ \tcp{Total data movement}
$B_{\text{bank}} \gets B_{\text{total}} / N_{\text{banks}}$

$n_{\text{rb}} \gets \lceil B_{\text{bank}} / R \rceil$ \tcp{Row buffer fills}
$t_{\text{fill}} \gets R / B$ \tcp{Time per fill}
$T_{\text{rb}} \gets n_{\text{rb}} \cdot t_{\text{fill}}$

$T_{\text{pe}} \gets B_{\text{bank}} / W \cdot \tau$ \tcp{PE compute time}

\eIf{asymmetric quant enabled}{
    $R_{\text{ops}} \gets (M+N) \cdot K$ \tcp{Reduction elements}
    $T_{\text{red}} \gets R_{\text{ops}}/N_{\text{banks}} \cdot e / W \cdot \tau$
    $T_{\text{pe}} \gets T_{\text{pe}} + T_{\text{red}}$
}{}

$T_{\text{latency}} \gets t_{\text{fill}} + \max((n_{\text{rb}}-1) \cdot t_{\text{fill}}, T_{\text{pe}})$
\end{algorithm}

Algorithm~\ref{alg:pim} captures the key insight that row-buffer fills and PE computation are pipelined: after filling the first row buffer, the PE begins computation while subsequent buffers fill from DRAM.
The total latency is the first fill time plus the maximum of remaining fill time and total PE time.

\subsection{Score@V on GPU}

After Q@K and softmax produce float16 attention scores $S \in \mathbb{R}^{M \times N}$, the final output is $O = S \cdot V$.
Since $S$ is float16 and $V$ must be dequantized from 2-bit, this operation is float16-intensive.
We route this to GPU, which offers massive float16 throughput (989 TFLOPS for B100) that far exceeds what near-bank PEs can provide.

For the dequantization of $V$, GPU reads 2-bit $V$ from memory, applies $v = s_V(\hat{v} - z_V)$, and computes the matrix multiply.
This is efficiently handled by GPU tensor cores.

\subsection{Heterogeneous Timeline}

\pimattn{} models GPU and PIM as independent execution engines with separate timelines, advancing concurrently:
\begin{equation}
T_{\text{total}} = \max(T_{\text{PIM}}^{\text{QK}}, T_{\text{GPU}}^{\text{softmax}} + T_{\text{GPU}}^{\text{SV}})
\end{equation}

In practice, since Q@K on PIM and softmax+Score@V on GPU are sequential (softmax depends on Q@K output), the total latency is the sum.
However, the \emph{per-layer} latency benefits from hiding PIM Q@K latency behind GPU projection operations (QKV projection, output projection), which execute concurrently on GPU while PIM processes attention.

% TODO Figure 5: Execution timeline — GPU (QKV proj) || PIM (Q@K) → GPU (softmax+SV+O proj), showing overlap

\subsection{KV Cache Quantization Overhead}

In decode mode, each newly generated token produces new key and value vectors that must be quantized before appending to the KV cache.
We model this quantization overhead on GPU:
\begin{equation}
T_{\text{quant}} = \max\left(\frac{2 \cdot d_{\text{kv}} \cdot B_{\text{elem}}}{\text{BW}_{\text{GPU}}}, \frac{4 \cdot d_{\text{kv}}}{\text{FLOPS}_{\text{GPU}}}\right)
\end{equation}
where $d_{\text{kv}} = n_{\text{kv\_heads}} \cdot d_{\text{head}}$ and the factor 4 accounts for 2 subtracts and 2 multiplies per element in asymmetric quantization.
For Llama-3-8B ($d_{\text{kv}}=1024$), this overhead is $<0.01\mu s$ per token, negligible relative to attention computation.
